\begin{table}[htbp]
\centering
\small
\caption{Complete experiment summary for the canonical-split and source-held-out studies. Historical Phase-5 results are excluded by construction.}
\label{tab:experimentsummary}
\begin{tabular}{lp{3.4cm}p{3.0cm}p{3.2cm}p{3.2cm}p{4.0cm}p{3.4cm}p{2.2cm}}
\toprule
Task & Scientific question & Models compared & Dataset protocol & Main result & Statistical result & Conclusion & Status \\
\midrule
6 & Which representation predicts ROP-Plus best on the canonical split? & A, B (RGB CNN), B embedding, C & canonical split 6203/1328/1331, single test pass & C AUC 0.925988 highest; A AUC 0.659582 lowest & descriptive; forms the reference for Task 7 & RGB-derived representations dominate scalar biomarkers & primary, complete \\
7 & Are the model differences statistically supported? & C vs B embedding & paired class-stratified bootstrap, 10,000 replicates, seed 42 & delta AUC +0.001085 & 95\% CI [-0.001691, +0.003868], p = 0.431 & five scalar biomarkers add no supported increment on the RGB embedding & primary, complete \\
8 & Does spatial vessel morphology add information? & D, E, F, G vs B & canonical split, single test pass & E AUC 0.932841; G AUC 0.934880 & E - B delta AUC +0.007938 (Task 8B closure) & supported spatial-vessel complementarity & complete \\
8B & Does the vessel finding survive a 10,000-replicate statistical closure? & E vs B embedding; G vs E & paired class-stratified bootstrap, 10,000 replicates, seed 42 & E - B AUC +0.007938; G - E AUC +0.002039 & E - B CI [+0.002630, +0.013131] p = 0.0031; Brier -0.027693 p = 0.0011; G - E CI [-0.000274, +0.004397] p = 0.090 & vessel increment supported; biomarker increment not & complete \\
9 & Can learned joint fusion beat concatenation? & H, I vs E, G & canonical split, staged CNN fine-tuning, single test pass & H AUC 0.925825; I AUC 0.916247 & H - E AUC -0.007016 (p = 0.131); I - H AUC -0.009578 (p = 0.0049, harmful) & learned interaction underperformed concatenation; biomarkers harmful in-network & secondary exploratory, complete \\
10 & Can biomarkers help by conditioning the vessel branch? & J1 (FiLM) vs J0 (matched control) & canonical split, frozen embeddings, head-only training & J0 AUC 0.929754; J1 AUC 0.927631 & J1 - J0 AUC -0.002123, p = 0.068; gamma within +/-0.3\% of 1; single-biomarker ablation changes probabilities by \textasciitilde{}2e-5 with 0 argmax flips & FiLM behaved near identity; no measurable conditioning effect & secondary exploratory, complete \\
11 & How well does this transfer to an unseen source? & B\_LOSO, E\_LOSO, G\_LOSO & three source-held-out folds; held-out source never used in fitting & FARFUM E 0.809779 vs B 0.787537; Farabi E 0.782898 vs B 0.769944 & E - B FARFUM +0.022242 p < 0.0001; Farabi +0.012954 p = 0.0018; G - E null in every fold & vessel benefit transfers; biomarker benefit does not; absolute performance drops by 0.12-0.15 AUC & diagnostic, complete \\
12 & Can class-conditional domain invariance improve unseen-source performance? & K1 (CC-DANN + MMD) vs K0 (matched control) & three source-held-out folds, target source fully unseen & K0 0.83293 vs K1 0.83013 (FARFUM); K0 0.79192 vs K1 0.79103 (Farabi) & K1 - K0 AUC -0.002796 (p = 0.143) and -0.000883 (p = 0.259); domain predictability unchanged at 0.88-0.98; representation shift reduced & outcome B: alignment reduced shift diagnostics but did not improve held-out disease performance & secondary exploratory, complete \\
13 & 1) Does a ROPDeepX-style dual-backbone RGB representation improve RGB performance? 2) Does the vessel representation remain complementary to it? & L, M0, M1 vs B embedding, E, G & canonical split, staged CNN training, single test pass & 1) No: M0 AUC 0.905978 vs B 0.924903. 2) Yes: M1 AUC 0.926913 vs M0 0.905978 & M0 - B -0.018925 (p = 0.0063); M1 - M0 +0.020936 (p < 0.0001); M1 vs E and M1 vs G null & vessel complementarity is architecture-robust; the tested dual-RGB representation was not better than the original embedding & secondary exploratory, complete \\
\bottomrule
\end{tabular}
\end{table}
